\documentclass[a4paper,11pt]{article}

% --- Encodage et Langue ---
\usepackage[utf8]{inputenc}
\usepackage[T1]{fontenc}
\usepackage[french]{babel}

% --- Mise en page et Typographie ---
\usepackage{geometry}
\geometry{margin=2.5cm, headheight=15pt}
\usepackage{microtype}
\usepackage{lmodern}

% --- Mathématiques ---
\usepackage{amsmath, amssymb}
\usepackage{centernot}

% --- Couleurs et Boîtes ---
\usepackage{xcolor}
\definecolor{primary}{RGB}{0, 85, 150}      % Bleu technique
\definecolor{secondary}{RGB}{240, 248, 255} % Bleu très clair pour les fonds
\definecolor{solcolor}{RGB}{245, 250, 245}  % Vert très clair pour les solutions
\definecolor{solframe}{RGB}{60, 120, 60}    % Vert foncé pour les bordures de solution

\usepackage{tcolorbox}
\tcbuselibrary{skins}

% --- Tableaux professionnels ---
\usepackage{booktabs}
\usepackage{array}

% --- Affichage de code / binaire ---
\usepackage{listings}
\lstset{
    basicstyle=\ttfamily\small,
    backgroundcolor=\color{gray!5},
    frame=single,
    framerule=0.5pt,
    rulecolor=\color{gray!50},
    numbers=none,
    breaklines=true,
    keepspaces=true,
    columns=fullflexible
}

% --- Titres ---
\usepackage{titlesec}
\titleformat{\section}
  {\color{primary}\Large\bfseries}
  {\thesection}{1em}{}[\titlerule]

% --- Gestion des solutions (Mettre \solutionsfalse pour cacher les solutions) ---
\newif\ifsolutions
\solutionstrue % <--- version prof
%\solutionsfalse % <--- version étudiant

\newtcolorbox{notebox}{
  colback=secondary,
  colframe=primary,
  fonttitle=\bfseries,
  title=Notes et Conseils,
  arc=4pt,
  boxrule=0.8pt
}

\newtcolorbox{solbox}{
  colback=solcolor,
  colframe=solframe,
  fonttitle=\bfseries,
  title=Solution,
  arc=4pt,
  boxrule=0.8pt
}

% --- Titre du document ---
\title{\textcolor{primary}{\Huge\bfseries TD 1 -- Électronique Numérique}\\[0.5em] \Large Représentation des nombres}
\author{}
\date{}

\begin{document}

\maketitle
\vspace{-2em}

\begin{notebox}
\begin{itemize}
    \item Le but de ce TD est de manipuler différentes représentations des \textit{nombres} et notamment de maîtriser les changements de base \textbf{sans calculatrice}.
    \item En dehors du TD, il est suggéré de vérifier que vous êtes suffisamment adroit pour les réaliser également \textit{avec} calculatrice ou à l'aide de votre langage de programmation préféré.
    \item \textit{Note : Tous les exercices ne seront pas nécessairement corrigés lors des séances de TD.}
\end{itemize}
\end{notebox}

\vspace{1em}

\section{Compter en base 2 et base 16}
\begin{enumerate}
    \item Compter de 0 à 15 en binaire, puis en hexadécimal. \textit{(On conservera ces valeurs dans un tableau pour la durée du TD)}

    \ifsolutions
    \begin{solbox}
    \begin{center}
      \renewcommand{\arraystretch}{1.2}
      \begin{tabular}{c c c}
        \toprule
        \textbf{Décimal} & \textbf{Binaire} & \textbf{Hexadécimal} \\
        \midrule
        0 & 0000 & 0 \\
        1 & 0001 & 1 \\
        2 & 0010 & 2 \\
        3 & 0011 & 3 \\
        4 & 0100 & 4 \\
        5 & 0101 & 5 \\
        6 & 0110 & 6 \\
        7 & 0111 & 7 \\
        8 & 1000 & 8 \\
        9 & 1001 & 9 \\
        10 & 1010 & A \\
        11 & 1011 & B \\
        12 & 1100 & C \\
        13 & 1101 & D \\
        14 & 1110 & E \\
        15 & 1111 & F \\
        \bottomrule
      \end{tabular}
    \end{center}
    \end{solbox}
    \fi

    \item Écrire les puissances de $2^n$ pour $n$ variant de 0 à 10. En déduire une valeur approchée de $\log_{10} 2$, puis $\log_{10} 5$. Avec un procédé similaire, trouver une valeur approchée de $\log_{10} 3$.

    \ifsolutions
    \begin{solbox}
    \begin{center}
      \renewcommand{\arraystretch}{1.2}
      \begin{tabular}{c c}
        \toprule
        $n$ & $2^n$ \\
        \midrule
        0 & 1 \\
        1 & 2 \\
        2 & 4 \\
        3 & 8 \\
        4 & 16 \\
        5 & 32 \\
        6 & 64 \\
        7 & 128 \\
        8 & 256 \\
        9 & 512 \\
        10 & 1024 \\
        \bottomrule
      \end{tabular}
    \end{center}
    \vspace{0.5em}
    On remarque que $2^{10} = 1024 \approx 10^3$. D'où :
    $$ \log_{10}(2^{10}) \approx \log_{10}(10^3) \implies 10 \log_{10} 2 \approx 3 \implies \log_{10} 2 \approx \frac{3}{10} = 0.3 $$
    (En réalité, $\log_{10} 2 \approx 0.301$ à $10^{-3}$ près. L'erreur est de 0.3\%).\\
    De même :
    $$ 5 = \frac{10}{2} \implies \log_{10} 5 = \log_{10} 10 - \log_{10} 2 \approx 1 - 0.3 = 0.7 $$
    Pour $\log_{10} 3$, on peut remarquer que $3^4 = 81 \approx 80 = 10 \times 2^3 \times 1$ (ou $3^2 \approx 10$). Une approximation courante utilise $3^4 \approx 10 \times 8 = 10 \times 2^3$ :
    $$ 4 \log_{10} 3 \approx 1 + 3 \log_{10} 2 \approx 1 + 3(0.3) = 1.9 \implies \log_{10} 3 \approx \frac{1.9}{4} = 0.475 $$
    (Valeur réelle $\approx 0.477$).
    \end{solbox}
    \fi

    \item Calculer la valeur décimale de $10101011_2$ en passant par les puissances de 2, puis en passant par les puissances de 16.

    \ifsolutions
    \begin{solbox}
    \begin{align*}
        10101011_2 &= 2^7 + 2^5 + 2^3 + 2^1 + 2^0 \\
                   &= 128 + 32 + 8 + 2 + 1 = \mathbf{171_{10}}
    \end{align*}
    En passant par l'hexadécimal (groupement par 4 bits : $1010$ et $1011$) :
    \begin{align*}
        1010\ 1011_2 &= \text{A}\text{B}_{16} \\
                     &= 10 \times 16^1 + 11 \times 16^0 \\
                     &= 160 + 11 = \mathbf{171_{10}}
    \end{align*}
    \end{solbox}
    \fi

    \item Écrire les puissances de $2^{-n}$ pour $n$ variant de 1 à 4.

    \ifsolutions
    \begin{solbox}
    \begin{center}
      \renewcommand{\arraystretch}{1.2}
      \begin{tabular}{c c}
        \toprule
        $n$ & $2^{-n}$ \\
        \midrule
        1 & 0,5 \\
        2 & 0,25 \\
        3 & 0,125 \\
        4 & 0,0625 \\
        \bottomrule
      \end{tabular}
    \end{center}
    \vspace{0.5em}
    Si on s'attribue 2 bits pour représenter (quantifier) une partie décimale, on a les équivalences suivantes :
    \begin{center}
      \begin{tabular}{c|c}
        \toprule
        Binaire & Décimal \\
        \midrule
        00 & 0,00 \\
        01 & 0,25 \\
        10 & 0,50 \\
        11 & 0,75 \\
        \bottomrule
      \end{tabular}
    \end{center}
    Pour 3 bits, on double la précision (pas de 0,125).
    \end{solbox}
    \fi
\end{enumerate}

\textit{Remarque :} Des nombres possédant une partie décimale sont transformés ici en de simples \textit{entiers}, en supposant que la position de la virgule est fixée à l'avance. Cette représentation est appelée \textbf{virgule fixe} (\textit{fixed-point}). Elle est très rapide à traiter par l'électronique, à coût matériel très faible, et est très utilisée en traitement du signal embarqué (DSP). À l'inverse, une représentation en \textbf{virgule flottante} est plus consommatrice de ressources (surface) et de temps, mais offre une plage dynamique et une précision beaucoup plus importantes (norme IEEE 754).

\section{Conversions entre bases}
\begin{enumerate}
    \item Convertir le nombre $39_{10}$ en binaire.
    \ifsolutions
    \begin{solbox}
    On effectue les divisions successives par 2 :
    \begin{align*}
        39 &= 2 \times 19 + 1 \\
        19 &= 2 \times 9 + 1 \\
        9 &= 2 \times 4 + 1 \\
        4 &= 2 \times 2 + 0 \\
        2 &= 2 \times 1 + 0 \\
        1 &= 2 \times 0 + 1
    \end{align*}
    En lisant les restes de bas en haut : $39_{10} = \mathbf{100111_2}$.
    \end{solbox}
    \fi

    \item Convertir le nombre $(0,375)_{10}$ en binaire.
    \ifsolutions
    \begin{solbox}
    On effectue les multiplications successives par 2 :
    \begin{align*}
        0,375 \times 2 &= 0,75 \quad (\text{partie entière } 0) \\
        0,75 \times 2  &= 1,50 \quad (\text{partie entière } 1) \\
        0,50 \times 2  &= 1,00 \quad (\text{partie entière } 1)
    \end{align*}
    En lisant les parties entières de haut en bas : $0,375_{10} = \mathbf{0,011_2}$. (On remarque que $0,375 = 3 \times 0,125 = 3 \times 2^{-3}$).
    \end{solbox}
    \fi

    \item En déduire la représentation binaire du nombre $(39,125)_{10}$.
    \ifsolutions
    \begin{solbox}
    On combine les résultats précédents (en notant que $0,125 = 0,001_2$) :
    $$ (39,125)_{10} = \mathbf{100111,001_2} $$
    \end{solbox}
    \fi

    \item Convertir le nombre $(0,2)_{10}$ en binaire. Que remarquez-vous ?
    \ifsolutions
    \begin{solbox}
    On effectue les produits successifs :
    \begin{align*}
        0,2 \times 2 &= 0,4 \quad (0) \\
        0,4 \times 2 &= 0,8 \quad (0) \\
        0,8 \times 2 &= 1,6 \quad (1) \\
        0,6 \times 2 &= 1,2 \quad (1) \\
        0,2 \times 2 &= 0,4 \quad (0) \quad \leftarrow \text{On retombe sur le début !}
    \end{align*}
    On trouve $0,2_{10} = \mathbf{0,00110011\dots_2} = 0,\overline{0011}_2$.
    \textit{Remarque :} Une fraction décimale finie peut avoir une représentation binaire infinie périodique, ce qui explique les erreurs d'arrondi en informatique.
    \end{solbox}
    \fi

    \item Convertir $(345,158)_{10}$ en octal.
    \ifsolutions
    \begin{solbox}
    \textbf{Partie entière (divisions successives par 8) :}
    \begin{align*}
        345 / 8 &= 43 \quad \text{reste } \mathbf{1} \\
        43 / 8  &= 5 \quad \text{reste } \mathbf{3} \\
        5 / 8   &= 0 \quad \text{reste } \mathbf{5}
    \end{align*}
    $\implies 345_{10} = 531_8$\\
    \textbf{Partie décimale (multiplications successives par 8) :}
    \begin{align*}
        0,158 \times 8 &= 1,264 \quad (\mathbf{1}) \\
        0,264 \times 8 &= 2,112 \quad (\mathbf{2}) \\
        0,112 \times 8 &= 0,896 \quad (\mathbf{0}) \\
        0,896 \times 8 &= 7,168 \quad (\mathbf{7})
    \end{align*}
    $\implies 0,158_{10} \approx 0,1207_8$.\\
    Résultat final : $\mathbf{531,1207_8}$.
    \end{solbox}
    \fi

    \item Convertir $(389)_{10}$ en hexadécimal.
    \ifsolutions
    \begin{solbox}
    Divisions successives par 16 :
    \begin{align*}
        389 / 16 &= 24 \quad \text{reste } \mathbf{5} \\
        24 / 16  &= 1 \quad \text{reste } \mathbf{8} \\
        1 / 16   &= 0 \quad \text{reste } \mathbf{1}
    \end{align*}
    Résultat : $\mathbf{185_{16}}$ (ou \texttt{0x185}).
    \end{solbox}
    \fi

    \item Convertir $011001111011101_2$ en hexadécimal.
    \ifsolutions
    \begin{solbox}
    On groupe les bits par paquets de 4 en partant de la droite (on peut ajouter un 0 à gauche) :
    $$ 0110\ 0111\ 1011\ 1101_2 $$
    $$ \Downarrow $$
    $$ 6\quad 7\quad \text{B}\quad \text{D}_{16} $$
    Résultat : $\mathbf{67BD_{16}}$. \textit{(Note : le corrigé original indiquait 33DD, ce qui correspondait à un regroupement erroné. Le regroupement correct donne 67BD)}.
    \end{solbox}
    \fi

    \item Convertir $(10110001101011,1111)_2$ en octal.
    \ifsolutions
    \begin{solbox}
    On groupe les bits par paquets de 3 en partant de la virgule :
    \begin{lstlisting}
      010 110 001 101 011 . 111 100
        2   6   1   5   3 .   7   4
    \end{lstlisting}
    Résultat : $\mathbf{26153,74_8}$.
    \end{solbox}
    \fi

    \item Convertir $(10110001101011,111101)_2$ en hexadécimal.
    \ifsolutions
    \begin{solbox}
    On groupe les bits par paquets de 4 :
    \begin{lstlisting}
      0010 1100 0110 1011 . 1111 0100
         2    C    6    B .    F    4
    \end{lstlisting}
    Résultat : $\mathbf{2C6B,F4_{16}}$.
    \end{solbox}
    \fi

    \item Le nombre $101110001_2$ est écrit en BCD (décimal codé binaire). Convertir le nombre décimal obtenu en binaire pur. Combien de bits sont alors nécessaires ? Conclure.
    \ifsolutions
    \begin{solbox}
    En BCD, on groupe par 4 bits pour lire les chiffres décimaux :
    \begin{lstlisting}
      0001  0111  0001
        1     7     1  => 171 en décimal
    \end{lstlisting}
    Conversion de $171_{10}$ en binaire pur :
    $171 = 128 + 32 + 8 + 2 + 1 = 2^7 + 2^5 + 2^3 + 2^1 + 2^0$
    $\implies \mathbf{10101011_2}$.\\
    Il faut \textbf{8 bits} en binaire pur, contre \textbf{12 bits} (3 $\times$ 4) en BCD.\\
    \textit{Conclusion :} Le BCD est moins compact et moins efficace pour le calcul, mais il est très utile pour l'affichage direct sur des afficheurs 7 segments, car chaque digit est isolé.
    \end{solbox}
    \fi
\end{enumerate}

\section{Nombres signés en complément à 2}
\subsection*{Rappels}
Jusqu'ici, nous nous sommes intéressés aux nombres positifs. Il existe plusieurs manières de représenter un nombre négatif $n$. La manière la plus intuitive serait de dédier un bit (le plus à gauche) pour préciser le signe (ex: $0$ pour $+$, $1$ pour $-$). Malheureusement, cette représentation possède deux défauts majeurs : il existe deux représentations du zéro ($+0$ et $-0$) et il faut modifier l'algorithme d'addition.

Les architectes des ordinateurs ont retenu une autre représentation plus pratique : le \textbf{complément à 2}. C'est cette représentation qui régit le fonctionnement des \textit{unités arithmétiques et logiques} (ALU) des ordinateurs.

Avec $k$ bits, on peut représenter l'ensemble des nombres entiers dans l'intervalle :
$$ [-2^{k-1},\ \dots,\ 2^{k-1}-1] $$
Le complément à 2 d'un entier $n$ est obtenu par le calcul de la quantité $2^k - n$. On peut vérifier aisément que :
\begin{itemize}
    \item Le complément à 2 du complément à 2 d'un entier $n$ est $n$ lui-même.
    \item Le complément à 2 de zéro est zéro (on néglige la retenue qui dépasse la taille $k$ fixée).
\end{itemize}

On doit retenir la \underline{formule de passage en complément à 2} d'un nombre négatif $-X$ :
\begin{equation}
    -X = \overline{X} + 1
\end{equation}
La barre figurant au-dessus du $X$ signifie ``inversion de tous les bits'' (complément à 1). Si l'entier $n$ est codé sur $k$ bits ($s_{k-1}\dots s_0$) en complément à 2, sa valeur décimale est :
\begin{equation}
    n_{10} = -s_{k-1}2^{k-1} + \sum_{i=0}^{k-2} s_i2^i
\end{equation}

\subsection*{Exercices}
\begin{enumerate}
    \item Déterminer la plage des valeurs des entiers signés sur $k=4$ bits. Représenter ces valeurs en binaire. Quelles sont les représentations des valeurs extrêmes ? de 0, 1, et -1 ?
    \ifsolutions
    \begin{solbox}
    Avec $k=4$ bits, la plage est $[-2^{4-1},\ 2^{4-1}-1] = \mathbf{[-8,\ 7]}$.
    \begin{center}
      \renewcommand{\arraystretch}{1.1}
      \begin{tabular}{c c | c c}
        \toprule
        Décimal & Binaire (Cà2) & Décimal & Binaire (Cà2) \\
        \midrule
        -8 & \textbf{1000} & 0 & 0000 \\
        -7 & 1001 & 1 & 0001 \\
        -6 & 1010 & 2 & 0010 \\
        -5 & 1011 & 3 & 0011 \\
        -4 & 1100 & 4 & 0100 \\
        -3 & 1101 & 5 & 0101 \\
        -2 & 1110 & 6 & 0110 \\
        -1 & 1111 & 7 & 0111 \\
        \bottomrule
      \end{tabular}
    \end{center}
    \textbf{Points clés :}
    \begin{itemize}
        \item Les nombres positifs sont codés ``normalement''.
        \item Alors qu'en non-signé on pouvait compter jusqu'à $+15$ avec 4 bits, ici on est limité à $+7$.
        \item Les nombres qui présentent un `1' en MSB (Most Significant Bit) sont forcément négatifs \underline{dans cette représentation}.
        \item Le codage de la valeur $-1$ est $1111_2$.
    \end{itemize}
    \end{solbox}
    \fi

    \item Quel est le complément à 2 du nombre $8$, codé sur $6$ bits ? Utiliser la formule (1) et vérifier le résultat en appliquant la formule (2).
    \ifsolutions
    \begin{solbox}
    \textbf{Formule (1) :} On inverse tous les bits et on ajoute 1.
    $8_{10} = 001000_2$ (sur 6 bits) \\
    Inversion : $110111_2$ \\
    Ajout de 1 : $110111_2 + 1_2 = \mathbf{111000_2}$.

    \textbf{Formule (2) :} Vérification de la valeur de $111000_2$ :
    $$ -1 \times 2^5 + 1 \times 2^4 + 1 \times 2^3 + 0 + 0 + 0 = -32 + 16 + 8 = \mathbf{-8} $$
    Le résultat est cohérent.
    \end{solbox}
    \fi

    \item Soustraire : $91_{10} - 108_{10}$ en passant par le complément à 2.
    \ifsolutions
    \begin{solbox}
    Attention ! 91 et 108 peuvent se représenter sur 7 bits, mais pas $-108$. Il faut donc passer à \textbf{8 bits} pour éviter les débordements.
    \begin{align*}
        91_{10} &= 64 + 16 + 8 + 2 + 1 = 01011011_2 \\
        108_{10} &= 64 + 32 + 8 + 4 = 01101100_2 \\
        -108_{10} &= \overline{01101100} + 1 = 10010011 + 1 = 10010100_2
    \end{align*}
    Addition en complément à 2 :
    \begin{lstlisting}
        01011011  ( 91)
      + 10010100  (-108)
      -----------
        11101111
    \end{lstlisting}
    Vérification du résultat $11101111_2$ :
    $$ -2^7 + (2^6 + 2^5 + 2^3 + 2^2 + 2^1 + 2^0) = -128 + (64 + 32 + 8 + 4 + 2 + 1) = -128 + 111 = \mathbf{-17} $$
    Le calcul est correct ($91 - 108 = -17$).
    \end{solbox}
    \fi

    \item Écrire les nombres suivants en binaire sur 6 bits : $+24, +31, -24, -31$. Passez ensuite à 7 bits. Que remarquez-vous ?
    \ifsolutions
    \begin{solbox}
    \textbf{Sur 6 bits :}
    \begin{itemize}
        \item $+24 = 011000$
        \item $+31 = 011111$
        \item $-24 = \overline{011000} + 1 = 100111 + 1 = 101000$
        \item $-31 = \overline{011111} + 1 = 100000 + 1 = 100001$
    \end{itemize}
    \textbf{Sur 7 bits :}
    \begin{itemize}
        \item $+24 = 0011000$
        \item $+31 = 0011111$
        \item $-24 = \overline{0011000} + 1 = 1100111 + 1 = 1101000$
        \item $-31 = \overline{0011111} + 1 = 1100000 + 1 = 1100001$
    \end{itemize}
    \textit{Remarque :} On observe que le bit de signe se \textbf{propage} vers la gauche (c'est l'extension de signe). Pour passer d'une taille à une autre, il suffit de dupliquer le MSB.
    \end{solbox}
    \fi

    \item Soustraire : $1100\ 0101\ 0100\ 0000_2 - 101000\ 1111\ 1111_2 - 110\ 1111\ 1110_2 - 10\ 1110_2$.
    \ifsolutions
    \begin{solbox}
    \textit{Piège classique :} Si on ne tient pas compte du format et qu'on complémente ``mal'' (sans aligner les tailles), on obtient un résultat faux. Il faut ramener tous les opérandes à la même taille (ici, 16 bits minimum, mais 17 bits pour être sûr avec des signés).

    Valeurs décimales (pour vérification) : $50496 - 10495 - 1790 - 46 = 38165$.

    En se ramenant à une représentation \textbf{signée sur 17 bits} (pour gérer correctement le bit de signe de l'opérande positif initial de 16 bits) :
    \begin{lstlisting}
       0 1100 0101 0100 0000  (+50496)
     + 1 1101 0111 0000 0001  (-10495 en C2 sur 17 bits)
     + 1 1111 1001 0000 0010  (-1790 en C2 sur 17 bits)
     + 1 1111 1111 1101 0010  (-46 en C2 sur 17 bits)
     -----------------------
     1 0 1001 0101 0001 0101
    \end{lstlisting}
    La retenue de 18ème bit est ignorée. Le résultat sur 17 bits est $0\ 1001\ 0101\ 0001\ 0101_2$, qui vaut bien $\mathbf{38165_{10}}$.

    \textbf{Règle générale :} Pour additionner/soustraire des nombres signés, il faut les coder sur $n$ bits, où :
    $$ n = 1 + \max \Big( \lfloor \log_2(|x_1|) \rfloor + 1,\ \dots,\ \lfloor \log_2(|x_m|) \rfloor + 1 \Big) $$
    Ici, $n = 1 + \max(16, 14, 11, 6) = 17$ bits.
    \end{solbox}
    \fi
\end{enumerate}

\section{Norme IEEE 754}
\paragraph{Rappel : norme IEEE 754 (simple précision, 32 bits)}
Un flottant est codé par trois champs : 1 bit de signe $s$, 8 bits d'exposant $e$ et 23 bits de mantisse $m$.
La valeur représentée est, pour la plupart des nombres :
\[
V = (-1)^s \times 2^{\,e-127} \times (1 + m)
\]
où $m = \sum_{i=1}^{23} b_i 2^{-i}$ est la partie fractionnaire. L'exposant est représenté en excédent (biais) à $127$. La mantisse est toujours normalisée avec un $1$ implicite (gain de place), sauf pour les cas particuliers suivants :
\begin{itemize}
    \item si $e=255$ et $m=0$, alors $V = \pm\infty$ (selon $s$) ;
    \item si $e=255$ et $m \neq 0$, alors $V = \mathtt{NaN}$ (Not a Number) ;
    \item si $e=0$ et $m=0$, alors $V = \pm 0$ ;
    \item si $e=0$ et $m \neq 0$, il s'agit d'un nombre dénormalisé, dont la valeur est $V = (-1)^s \times 2^{-126} \times m$.
\end{itemize}

\paragraph{Exercice :} Représenter le nombre $238,4375$ dans sa représentation IEEE 754 (32 bits) en virgule flottante.

\ifsolutions
\begin{solbox}
\textbf{1. Conversion en binaire :}\\
Partie entière : $238_{10} = 128 + 64 + 32 + 8 + 4 + 2 = \mathbf{11101110_2}$\\
Partie décimale : $0,4375 = 0,25 + 0,125 + 0,0625 = 2^{-2} + 2^{-3} + 2^{-4} = \mathbf{0,0111_2}$\\
Donc : $238,4375_{10} = \mathbf{11101110,0111_2}$.

\textbf{2. Normalisation :}\\
On déplace la virgule pour obtenir un seul `1` avant la virgule :
$$ 11101110,0111_2 = 1,11011100111_2 \times 2^7 $$
L'exposant réel est $n = 7$.

\textbf{3. Codage des champs :}
\begin{itemize}
    \item \textbf{Signe ($s$) :} Le nombre est positif, donc $s = \mathbf{0}$.
    \item \textbf{Exposant ($e$) :} $e = n + 127 = 7 + 127 = 134_{10}$.\\
    $134_{10} = 128 + 4 + 2 = \mathbf{10000110_2}$.
    \item \textbf{Mantisse ($m$) :} On prend la partie fractionnaire après la virgule ($11011100111$) et on complète avec des zéros à droite pour atteindre 23 bits :\\
    $\mathbf{110\ 1110\ 0111\ 0000\ 0000\ 000_2}$.
\end{itemize}

\textbf{4. Assemblage final (32 bits) :}
\begin{lstlisting}
  S |    Exposant (8)   |          Mantisse (23)           |
  0 | 1 0 0 0 0 1 1 0 | 1 1 0 1 1 1 0 0 1 1 1 0 0 0 0 0 0 0 0 0 0 0 0 |
\end{lstlisting}
Soit en hexadécimal : $\mathbf{0x436E7000}$.
\end{solbox}
\fi

\end{document}
